\section{Entering the Watch}\label{sec:entering}

Advent begins in the gathering dark. On the Saturday evening before the
First Sunday of Advent, at First Vespers, the Church lays aside green
vestments, puts on violet, and opens her new year by turning toward the
coming of the Lord. While late November and December grow loud with hurried
errands and premature feasting, the Roman liturgy keeps four weeks of
watchful quiet, shaped by the cry of the prophets, the preaching of John
the Baptist in the wilderness, and the silence of Mary awaiting her Son.

If you are opening Advent before First Vespers arrives, take a quiet
half-hour to settle a fixed time and an uncluttered corner for daily prayer
(Section~\ref{sec:prayer-home}). Choose a modest rule that can survive
ordinary workdays and family duties rather than an ambitious schedule that
breaks down by Wednesday. Spend those early evenings reflecting on the
threefold coming of Christ, how the Church learned to keep this vigil, and
how prayer, mercy, and confession belong together, so that when the priest
chants the opening antiphon of the First Sunday, you know what you are
keeping watch for.

If you are arriving late---halfway through December, on a crowded weekday
evening, tempted to think the season has already passed you by---do not
wait for next Sunday or next year. Turn straight to today's date in
Appendix~\ref{app:daily-course}, read the appointed Scripture and patristic
commentary, and pray the day's prayer before sleep. In the Gospel of
Matthew, the householder goes out to the marketplace even at the eleventh
hour and asks, ``Why stand you here all the day idle?'' Hearing that no man
has hired them, he sends them straight into the field: ``Go ye also into my
vineyard'' (Matt 20:6--7). At evening, the Master of the vineyard pays the
latecomer the full wage given to those who bore the heat of the day: ``I
will also give to this last even as to thee'' (Matt 20:14).

Keeping Advent in the 1962 Roman Rite draws you into a living order of
worship whose terms quickly grow familiar. At the center stands the
\textbf{Mass}. Each Sunday of Advent, together with the three Winter Ember
Days of the third week and the Vigil of the Nativity on 24~December, has
its own proper chants, readings, and prayers. Listen closely to the
\textbf{Collect} (\textit{Oratio}) before the Epistle, where the priest
gathers the longing of the whole assembly into a single sentence. On
ordinary weekdays outside the Ember week, when no saint's feast takes
precedence, the Mass repeats the preceding Sunday's prayers and readings
without the Alleluia, letting the sober cry of the week sink in.

Around the Mass turns the \textbf{Divine Office}, the Church's daily round
of psalms, hymns, and readings: \textbf{Matins} in the night watches, with
its readings from the prophet Isaiah and the Fathers; \textbf{Lauds} at
dawn; the shorter daytime hours of \textbf{Prime}, \textbf{Terce},
\textbf{Sext}, and \textbf{None}; \textbf{Vespers} at dusk, with the hymn
\textit{Creator alme siderum}, ``Creator of the starry pole,'' and Mary's
\textit{Magnificat}; and \textbf{Compline} before sleep, ending with the
Marian antiphon \textit{Alma Redemptoris Mater}, begging the Mother of the
Redeemer to aid a fallen people striving to rise. Lay Catholics are not
bound to the full clerical Office, yet praying even one hour---or saying
the season's collects, antiphons, and hymns gathered in
Appendix~\ref{app:prayers}---joins the quiet of a household to the prayer
of the universal Church.

At the head of each Sunday, Ember, and Vigil Mass in the Missal stands a
\textbf{station} (\textit{Statio}), naming one of the ancient basilicas of
Rome where the pope and his people assembled for that liturgy: St.~Mary
Major, Holy Cross in Jerusalem, St.~Peter's, and the Twelve Holy Apostles.
Those names tie every parish and home to the apostolic city and to the
saints who kept this watch before us.

Psalms follow the Latin liturgical numbering first, with the Hebrew number
in brackets where they differ---such as Psalm~24~[25], whose cry \textit{Ad
te levavi}, ``To thee, O Lord, have I lifted up my soul,'' opens the First
Sunday of Advent. Details on Scripture translations and liturgical editions
appear in Appendix~\ref{app:about}.
